
\documentclass[11pt]{article}
\usepackage{geometry}                
\geometry{letterpaper}                   

\usepackage{graphicx}
\usepackage{amssymb}
\usepackage{epstopdf}

\usepackage[utf8]{inputenc}
\usepackage[english]{babel}
\usepackage{natbib}
\bibliographystyle{unsrtnat}
\renewcommand\bibsection{\section{\refname}}

\usepackage{amssymb, amsmath}
\usepackage[procnames]{listings}
\usepackage{color}
\usepackage[english]{varioref}

\DeclareGraphicsRule{.tif}{png}{.png}{`convert #1 `dirname #1`/`basename #1 .tif`.png}

\title{Vector based navigation of desert ants}
\author{Josua Graf, Noah Zarro}
\date{date} 

\graphicspath{ {./images/} }

\definecolor{keywords}{RGB}{255,0,90}
\definecolor{comments}{RGB}{0,0,113}
\definecolor{red}{RGB}{160,0,0}
\definecolor{green}{RGB}{0,150,0}
\lstset{language=Python, 
        basicstyle=\ttfamily\small, 
        keywordstyle=\color{keywords},
        commentstyle=\color{comments},
        stringstyle=\color{red},
        showstringspaces=false,
        identifierstyle=\color{green},
        procnamekeys={def,class}}

\begin{document}



\thispagestyle{empty}

\begin{center}
\includegraphics[width=5cm]{ETHlogo.eps}

\bigskip


\bigskip


\bigskip


\LARGE{ 	Lecture with Computer Exercises:\\ }
\LARGE{ Modelling and Simulating Social Systems\\}

\bigskip

\bigskip

\small{Project Report}\\

\bigskip

\bigskip

\bigskip

\bigskip


\begin{tabular}{|c|}
\hline
\\
\textbf{\LARGE{Vector based navigation of desert ants}}\\
\textbf{\LARGE{}}\\
\\
\hline
\end{tabular}
\bigskip

\bigskip

\bigskip

\LARGE{Josua Graf, Noah Zarro}



\bigskip

\bigskip

\bigskip

\bigskip

\bigskip

\bigskip

\bigskip

\bigskip

Zurich\\
\today
%Dec 2018\\

\end{center}



\newpage

%%%%%%%%%%%%%%%%%%%%%%%%%%%%%%%%%%%%%%%%%%%%%%%%%

\newpage
\section*{Agreement for free-download}
\bigskip


\bigskip


\large We hereby agree to make our source code for this project freely available for download from the web pages of COSS. Furthermore, we assure that all source code is written by ourselves and is not violating any copyright restrictions.

\begin{center}

\bigskip


\bigskip


\begin{tabular}{@{}p{3.3cm}@{}p{6cm}@{}@{}p{6cm}@{}}
\begin{minipage}{3cm}

\end{minipage}
&
\begin{minipage}{6cm}
\vspace{2mm} \large Josua Graf

 \vspace{\baselineskip}

\end{minipage}
&
\begin{minipage}{6cm}

\large Noah Zarro

\end{minipage}
\end{tabular}


\end{center}
\newpage

%%%%%%%%%%%%%%%%%%%%%%%%%%%%%%%%%%%%%%%



% IMPORTANT
% you MUST include the ETH declaration of originality here; it is available for download on the course website or at http://www.ethz.ch/faculty/exams/plagiarism/index_EN; it can be printed as pdf and should be filled out in handwriting


%%%%%%%%%% Table of content %%%%%%%%%%%%%%%%%

\tableofcontents

\newpage

%%%%%%%%%%%%%%%%%%%%%%%%%%%%%%%%%%%%%%%



\section{Abstract}

\section{Individual contributions}

\section{Introduction and Motivations}
	\subsection{Motivation}
		we wanted to do an unconventional poject, simulating emergency situations has been done multiple times
	\subsection{Base Paper}
		description of base paper (who did it, roughly what was it all about)
\section{Description of the Model}
	\subsection{Behaviour of the ant}
		how did Wehner think global and local vector worked (\cite{wehner})
	\subsection{Conducted experiments}
		description of the experiments Wehner conducted (explaination of channel, cylinder, etc.)
\section{Implementation}
	with code examples
	\subsection{Main Concept}
		we use an iterative aproach, calculate local and global vector new in every step, merge them together(explain how). Save steps, plot them
	\subsection{Global Vector}
		explain how it gets calculated, and randomized
	\subsection{Local Vector}
		explain how it gets calculated and influenced by surroundings
\section{Simulation Results and Discussion}
	\subsection{Simulation Results}
		what did we get, is it the same
	\subsection{Discussion}
		what was different, probably the randomization was no as in real life
\section{Summary and Outlook}
	\subsection{Summary}
		it worked pretty well, Wehner did a good job
	\subsection{Outlook}
		we probably will not continue our research about desert ants, maybe Wehner will
%\section{References}
	\bibliography{bibliography}
\section{Latex-Stuff}

Beispielverweisung:

Hier wird auf dieses Bild verwiesen (fig.~\vref{fig:examplefigure})

Python Code:
\begin{lstlisting}[caption=Python code example]
import numpy as np

# show numbers 0-9
print ('Numbers from 0-9')
for i in range(10):
	print(i)
\end{lstlisting}


\begin{figure}[h!]
	\includegraphics[width=0.5\textwidth]{example.png}
	\caption{Beispielbild, in Ordner images}
	\label{fig:examplefigure}
\end{figure}

Beispielliste:

\begin{itemize}
	\item erster Eintrag
	\item zweiter Eintrag
\end{itemize}




\end{document}  

\printbibliography

 
